\section{A descriptive coding-agent case study}
\label{sec:coding-case-study}

\subsection{Design, data and evidence status}
We examine an archived coding-agent experiment to assess which parts of the evaluation contract were implemented
and which scientific questions remain unresolved. The environment uses 427 sanitized MBPP tasks and 164 HumanEval
tasks \citep{austin2021program,chen2021code}. These are public function-synthesis benchmarks, not a sample of all
agent tasks. A frozen split assigns 30 tasks to the pilot, 231 to training and 330 to evaluation; the evaluation
partition contains 239 MBPP and 91 HumanEval tasks. The intended primary inference is conditional on these fixed
benchmark tasks. This section reports observed comparisons and retrospective diagnostics, not independently
validated coverage or a completed confirmatory test of the final research claims.

The two configurations are Qwen2.5-3B-Instruct and Qwen2.5-7B-Instruct, both Q4\_K\_M, with decoding temperature
0.7, top-$p$ 0.95 and a 1,024-token output cap per call. We denote them small and large. The first model call
generates a candidate function. A second or third call is eligible only after failure of visible validation;
passing that validation causes immediate submission. Hidden benchmark tests score the submitted candidate but
do not govern deployed stopping. The visible checks were generated and certified against reference solutions
before collection, with duplicate hidden-test inputs removed. Configurations, task partitions, routing draws,
visible checks and policy definitions are frozen in the repository protocol and manifests.

The randomized log contains 4,488 episodes, eight per training or evaluation task. Within each task, the initial
assignment is blocked to four small and four large calls; later eligible assignments have probability $1/2$ per
model. The evaluation log contains 2,640 episodes. A policy learned on training tasks is frozen before the
3,960 fresh target-policy executions, comprising six policies with two runs on each of the same 330 evaluation
tasks. These repeated runs are not 3,960 independent tasks. The archived live outcomes and utility identities
have been independently reconstructed; the independent review did not rerun models or hidden-test verification.

\subsection{Observed success and resource use}
We retain the original utility $U=Y-0.01N_S-0.03N_L$, where $Y$ is hidden-test success and $N_S,N_L$ count model
calls. These unitless penalties are a declared preference, not measured money or energy. Table~\ref{tab:coding-live}
reports all six evaluated policies, without selecting a favorable subset. First-failure escalation starts small
and uses large at subsequent eligible calls. Class-tailored starts small and uses large after an assertion failure,
small after other failures. Soft escalation selects large with probability $1/2$ initially and $2/3$ subsequently. The learned
policy is discussed below. All policies share the same stopping rule and maximum of three model calls.

\begin{table}[htbp]
\centering
\small
\begin{tabular}{lrrrr}
\toprule
Policy & Success & Utility & Large calls & All calls \\
\midrule
Always small & .6106 & .5956 & .0000 & 1.5000 \\
Always large & .7167 & .6777 & 1.3000 & 1.3000 \\
First-failure escalation & .6470 & .6245 & .4152 & 1.4152 \\
Class-tailored & .6515 & .6319 & .2742 & 1.4152 \\
Soft escalation ($\delta=2$) & .6758 & .6464 & .7818 & 1.3758 \\
Learned fixed repair schedule & .7091 & .6727 & 1.1545 & 1.3333 \\
\bottomrule
\end{tabular}
\caption{Observed evaluation-cohort means, 660 episodes on 330 tasks per policy. Call counts are per episode.
These are descriptive success/resource comparisons, not an optimized frontier or a demonstration of equal quality,
policy superiority or deployment savings.}
\label{tab:coding-live}
\end{table}

The learned policy has five fewer successes than always-large among 660 executions per policy, a difference of
$-0.007576$ in success and $-0.005000$ in utility. It uses 96 fewer large calls but 118 more small calls,
or 22 more total calls. These realized differences do not establish equivalence or a positive advantage.
For the two frozen policies, revaluing the observed outcomes at penalties $c_S,c_L$ gives
\[
\widehat\Delta_U(c_S,c_L)=\frac{-5+96c_L-118c_S}{660}.
\]
The empirical crossing is $c_L=0.064375$ if $c_S=0.01$ remains fixed. If both original penalties are scaled
together, the crossing multiplier is $50/17$, giving $c_L\simeq0.088235$ and $c_S\simeq0.029412$.
This retrospective sensitivity neither changes the original endpoint nor identifies what a newly trained policy
would do at different penalties. A separate training-only refit without call penalties leaves the reachable
learned schedule unchanged.

\subsection{What the learned policy and stopping rule permit}
Independent backward-regression reconstruction reproduces all 17 frozen policy entries. On reachable states,
however, the deployed rule is simply large, then small after a failure, then large after a further failure.
All 660 live paths follow this fixed schedule at eligible stages: 542 stop after large, 16 after large/small and
102 reach large/small/large. Off-path differences in the policy table do not establish history-dependent action
selection on the policy's own trajectories. Its initial key uses only benchmark identity, preventing initial
routing by task difficulty within either benchmark. An improvement over always-small would therefore not isolate
the value of adapting to the evolving failure history.

Feedback quality further restricts the available decisions. After certification, 60 of the 330 evaluation tasks
have no surviving visible checks. On those tasks, all 120 learned and 120 always-large executions stop after
the first call, although respectively 61 and 65 fail hidden tests. Across the learned cohort, 104 first-call
submissions pass visible validation but fail hidden tests, accounting for 54.2\% of its 192 final failures.
Only 118 of 660 episodes reach a first repair. These observations identify an eligibility limitation, not the
counterfactual claim that additional calls would correct those failures. A different feedback or stopping rule
would define a new harness and require a prospective evaluation; hidden verification cannot become a deployed
stopping signal. The present data are retained in full, including zero-check tasks.

\subsection{Offline/live discrepancies and their unresolved interpretation}
Table~\ref{tab:coding-calibration} retains every live policy in the offline/live comparison. Offline entries are
the archived three-fold task-split DR estimates; the class-tailored estimate and its original fitting procedure
were independently reconstructed, while the other DR entries are transcribed from the pinned analysis summaries.
This distinction matters: a common table is not evidence that every estimator was independently reimplemented.

\begin{table}[htbp]
\centering
\small
\begin{tabular}{lrrr}
\toprule
Policy & Offline DR & Fresh live & DR minus live \\
\midrule
Always small & .5946 & .6106 & $-.0160$ \\
Always large & .7114 & .7167 & $-.0052$ \\
First-failure escalation & .6306 & .6470 & $-.0164$ \\
Class-tailored & .6149 & .6515 & $-.0366$ \\
Soft escalation ($\delta=2$) & .6689 & .6758 & $-.0069$ \\
Learned fixed repair schedule & .7260 & .7091 & $+.0169$ \\
\bottomrule
\end{tabular}
\caption{Hidden-test success estimates on the same 330 evaluation tasks. Live means have execution variability;
these differences are not bias estimates against known truth. No validated-coverage claim or ranking of estimator
accuracy follows from this single dataset.}
\label{tab:coding-calibration}
\end{table}

Class-tailored IPW success is .610606, DR is .614871 and live success is .651515. Thus the DR adjustment narrows
the discrepancy with live outcomes. Independent checks reproduce the recorded policy actions, assignment draws,
probabilities and calculations; they find no mismatch in those inspected components. A first-candidate diagnostic,
before adaptive repair, already differs between initial-small log and class-tailored live executions (.565152
versus .590909). This localization is retrospective and does not identify the cause. Finite execution variation,
differences in the execution mechanism and shortcomings in inference remain distinct possibilities. Selecting a
new estimator because it agrees with these live means would not resolve the calibration question.

The prior common-five replay comparison has mean absolute discrepancies .016061 for donor replay and .018235
for DR against noisy live means. It does not demonstrate DR superiority. Four of those deterministic targets
reduce to fixed eligible-stage schedules, and donor fallback remains present. The exact adaptive replay
counterexample in Section~\ref{adaptive-donor-replay} establishes a possible failure mechanism, not its empirical
importance in this particular dataset.

\subsection{Branch target, remaining inference and implications}
The archived branch experiment samples 200 of 564 logger-reached first-failure prefixes, spanning 103 sampled
tasks, and records two fresh continuations per model per prefix: 800 continuations in total. Each assigned model
is used throughout the remaining eligible stages. The branch contrast is .120000 and the pooled log contrast
is .134654 for stay-large minus stay-small. This is a local continuation comparison on the logger's prefix
distribution, not the learned small/large continuation, class-tailored behavior, or a whole-policy value.
Recovery after an archive-loss incident further requires an observation-process justification: retained records
and reconstructed transcript bytes do not prove that lost outcomes or execution state were exchangeable.

The primary fixed-benchmark source/selection/execution inference is not yet validated. Conditioning on the
entire observed prefix frame yields the conditional result in Section~\ref{sec:fixed-frame-branches}, but changes
the conditioning set and treats the realized log comparator as fixed. The separate iid source-population result
does not automatically apply to tasks with zero eligible prefixes or zero arm contributions. Neither the existing
derivative scale nor a variance estimate alone supplies the missing coverage argument. We consequently report
the branch points without presenting the exploratory band as a validated confidence interval.

The case study supports a narrower contribution than demonstrated adaptive improvement: a prospective routing
archive that makes design and interpretation failures inspectable. Informative deployable feedback, competitive
fixed-schedule and learned-router controls, an explicit success/resource criterion, and inference for the declared
sampling design remain necessary. Future choices will be developed outside this evaluation cohort and tested
on fresh data. A null or unfavorable comparison is retained as evidence; it is not repaired by changing the
endpoint or treating noisy estimated subgroup gains as an oracle upper bound.
